\chapter*{پیش‌گفتار مترجم}
\addcontentsline{toc}{chapter}{پیش‌گفتار مترجم}
\markboth{پیش‌گفتار مترجم}{پیش‌گفتار مترجم}

فرآیند ترجمه‌ی کتاب «خط فرمان لینوکس» (\en{The Linux Command Line}) اثر ویلیام شاتس، با اخذ اجازه‌ی کتبی از ایشان به‌عنوان نویسنده و صاحب‌حقوق اثر آغاز شد. این اثر که در انتشارات \en{No Starch Press} به چاپ رسیده، تحت مجوز \en{Creative Commons} (\en{BY-NC-ND 3.0}) به‌صورت آزاد عرضه شده و تاکنون به زبان‌های متعددی برگردانده شده است. نسخه‌ی پیش‌رو بر اساس ویرایش ششم نسخه‌ی اینترنتی (\en{Internet Edition}) تنظیم شده و با هدف ارائه‌ی منبعی دقیق و منسجم برای علاقه‌مندان فارسی‌زبان به این حوزه‌ی تخصصی تدوین گشته است.

انجام تمام مراحل این پروژه، از ترجمه تا نگارش و صفحه‌آرایی، به تنهایی صورت گرفت و با مسائل فنی متعددی همراه بود. یکی از بخش‌های چالش‌برانگیز کار، مدیریت تداخل ساختاری میان متن فارسی (راست‌به‌چپ) و عناصر فنی انگلیسی (چپ‌به‌راست) بود. حفظ یکپارچگی پاراگراف‌ها و ترتیب صحیح نویسه‌ها در هنگام درج دستورات، مستلزم بازبینی‌های مکرر و دقیق در تمام مراحل حروف‌چینی بود تا خوانش متن دچار اختلال نشود.

در حوزه‌ی واژه‌گزینی نیز با توجه به سرعت تغییرات تکنولوژی و چالش‌های رایج در انتخاب معادل‌های فنی، رویکردی متعادل اتخاذ شد. در سرتاسر متن، معادل‌های فارسی در کنار اصطلاحات اصلی (به زبان انگلیسی) آورده شده‌اند.

در طی بازگردانی متن، تلاش شد تا با وسواس در انتخاب واژگان و ساختار جملات، متنی روان و دقیق حاصل شود. پس از اتمام ترجمه، بازخوانی چندباره‌ی کل اثر انجام شد و دقت نظر ویژه‌ای اعمال شد تا استانداردسازی لحن، یکدست‌سازی ساختارها و درج صحیح علائم نگارشی به‌درستی صورت گیرد.

این اثر، نخستین ویرایش از این ترجمه است و بی‌تردید قابلیت ارتقا دارد. بدین‌وسیله از خوانندگان، متخصصان و فعالان حوزه‌ی لینوکس تقاضا دارم هرگونه ایراد احتمالی در ترجمه یا پیشنهاد برای بهبود معادل‌های فنی را از طریق نشانی پست الکترونیک \lr{MohammadMehdinejad@Proton.Me} با اینجانب در میان بگذارند تا در ویرایش‌های بعدی لحاظ شود.

امید است این کتاب گامی هرچند کوچک در جهت تسهیل یادگیری خط فرمان لینوکس برای جامعه‌ی فارسی‌زبان باشد.

\vspace{1.5em}
\begin{flushleft}
خرداد ۱۴۰۴ – محمد مهدی‌نژاد
\end{flushleft}